\documentclass[12pt]{article}
\usepackage[T1]{fontenc}
\usepackage[utf8]{inputenc}
\usepackage[brazil]{babel}
\usepackage[margin=1in]{geometry}
\setlength{\parindent}{0pt}
\begin{document}
\section*{Tema 55}
Processo: PEDILEF 0355079-05.2005.4.03.6301/ SP\\
Situacao: Julgado\\
Ramo: DIREITO CIVIL\\
Questao: Saber se é possível cumulação de juros moratórios com juros remuneratórios progressivos e expurgos inflacionários incidentes sobre saldo de conta vinculado ao FGTS, bem como qual a taxa de juros de mora aplicável.\\
Tese: São devidos, além dos juros progressivos sobre os saldos fundiários, juros moratórios, previstos no art. 406 do Código Civil e art. 161, §1º do Código Tributário Nacional, pela taxa SELIC, contados a partir da citação até a data do pagamento. Vide Tema 176 do STJ - Recursos Repetitivos.\\
Fonte: https://www.cjf.jus.br/phpdoc/virtus/
\end{document}
